\documentclass[10pt]{book}

% 日本語（LuaLaTeX）
	\usepackage{luatexja}
	\usepackage{luatexja-fontspec}

% 色
	\usepackage{xcolor}

% コード
	\usepackage{listings}
	\usepackage{tcolorbox}
	\tcbuselibrary{skins, breakable, listings}

% 見出し
	\usepackage{titlesec}

% 行間
	\usepackage{setspace}

% 余白
	% \usepackage{geometry}
	% layout.tex で geometry のオプション指定があるので、ここでは読み込まない

% ヘッダ・フッタ
	\usepackage{fancyhdr}

% 図表
	\usepackage{graphicx}
	\usepackage{caption}

% TikZ（章扉）
	\usepackage{tikz}

% 目次
	\usepackage{tocloft}

% コラムアイコン
	\usepackage{fontawesome5}

% listing の位置調整
	\usepackage{placeins}
	\usepackage{float}

% ---------------------------------------------------------
% 日本語フォント設定
% ---------------------------------------------------------
% 日本語フォント（Windows）
\setmainfont{Noto Serif}      % 欧文本文：セリフ
\setsansfont{Noto Sans}       % 欧文ゴシック
\setmainjfont{Noto Serif JP}  % 和文明朝（本文）
\setsansjfont{Noto Sans JP}   % 和文ゴシック（見出し用）

% 日本語の太字を有効化（luatexja）
\ltjsetparameter{jacharrange={-2}} % 日本語の太字を許可
% 日本語ゴシック体（キャプション専用）
\newjfontfamily\captionjfont{Noto Sans JP}[
  UprightFeatures = { Weight = 400 }, % Regular
  BoldFeatures    = { Weight = 700 }  % Bold
]
\newjfontfamily\JPSans{Noto Sans JP}[
  UprightFeatures = { Weight = 400 }, % Regular
  BoldFeatures    = { Weight = 700 }  % Bold
]
% コード用日本語ゴシック体
\newjfontfamily\codejfont{Noto Sans JP}[
  UprightFeatures = { Weight = 400 },
  BoldFeatures    = { Weight = 700 }
]
\definecolor{chapterblue}{HTML}{DCE4F2}

\definecolor{sectionblue}{HTML}{c9d1e8} %よりグレーに近い青に変更（6F8FBC → c9d1e8）
%\definecolor{sectionblue}{HTML}{4A6FA5}

\definecolor{framegray}{HTML}{D0D7DE}

\setstretch{1.4}

\usepackage[
  paperwidth=148mm,
  paperheight=210mm,
  inner=20mm,
  outer=10mm,
  top=18mm,
  bottom=20mm,
  heightrounded,
  twoside
]{geometry}

% ---------------------------------------------------------
% ヘッダ・フッタ（落ち着いた市販書籍風）
% ---------------------------------------------------------
\pagestyle{fancy}

% --- ヘッダのフォント（ゴシック体） ---
\newcommand{\HeaderFont}{\small\sffamily\JPSans}

% --- 章タイトル（chapter）を記録 ---
\renewcommand{\chaptermark}[1]{%
  \markboth{\thechapter\quad #1}{}%
}


% --- 節タイトル（section）を記録 ---
\renewcommand{\sectionmark}[1]{%
  \markright{\thesection\quad #1}%
}


% --- ヘッダ設定 ---
\fancyhf{} % 既存の設定をクリア

% 偶数ページ（左ページ）：章タイトル
\fancyhead[LE]{\HeaderFont \leftmark}

% 奇数ページ（右ページ）：節タイトル
\fancyhead[RO]{\HeaderFont \rightmark}

% ページ番号を外側寄せ
\fancyfoot[LE]{\HeaderFont \thepage} % 偶数ページ（Left Even）
\fancyfoot[RO]{\HeaderFont \thepage} % 奇数ページ（Right Odd）

% ヘッダの下線を消す（技評風）
\renewcommand{\headrulewidth}{0pt}

% --- plain スタイル（目次・章扉などで使われる）を外側寄せに再定義 ---
\fancypagestyle{plain}{
  \fancyhf{}
  % ページ番号を外側寄せ
  \fancyfoot[LE]{\HeaderFont \thepage} % 偶数ページ（左）
  \fancyfoot[RO]{\HeaderFont \thepage} % 奇数ページ（右）
  \renewcommand{\headrulewidth}{0pt}   % ヘッダ線なし（技評風）
}


% 章扉背景（ページ全面を確実に塗りつぶす）
\newcommand{\ChapterBackground}{%
  \begin{tikzpicture}[remember picture,overlay]
    % ページ全体より少し大きめに塗る（塗り残し防止）
    \fill[chapterblue]
      ([xshift=-5mm,yshift=5mm]current page.north west)
        rectangle
      ([xshift=5mm,yshift=-5mm]current page.south east);
  \end{tikzpicture}%
}

% 章デザイン（bbe風：背景色＋装飾）
\titleformat{\chapter}[display]
  {\bfseries\Huge}
  {%
    \ChapterBackground
    \vspace*{3cm}
    \filright\Large\textsf{CHAPTER\ \thechapter}
  }
  {1ex}
  {%
    \titlerule[1.2pt]
    \vspace{2ex}
    \filright\Huge
  }
  [\vspace{2cm}]
%\newcommand{\coverpage}[1]{%
  \clearpage
  \thispagestyle{empty}
  \begin{center}
    \vspace*{0pt}
    \includegraphics[width=\paperwidth]{#1}
  \end{center}
  \clearpage
}

% --------------------------------------------------------

% セクション見出しデザイン
% --------------------------------------------------------
% 見出しデザイン（技術書風）


\makeatletter
\def\SectionTitle{}
\def\ttl@save@title#1{\gdef\SectionTitle{#1}}
\makeatother

% --- Section：青背景＋黒文字＋ゴシック太字 ---
\let\oldsection\section

\renewcommand{\section}[1]{%
  \refstepcounter{section}%
  \SectionHeading{\thesection}{#1}%
}

\newcommand{\SectionHeading}[2]{%
  \sectionmark{#2} % ← ★ これが右ページの柱に必要
  \addcontentsline{toc}{section}{\protect\numberline{\thesection}#2}
  \vspace{1.5em}%
  \noindent
  \colorbox{sectionblue}{% sectionblueから変更
    \parbox{\linewidth}{%
      \vspace{2pt}%
      {%
        \Large\bfseries\sffamily\color{black}%
        \JPSans % ← 和文だけ Noto Sans JP に切り替える
        #1\quad #2%
      }%
      \vspace{2pt}%
    }%
  }%
  \vspace{1.2em}%
}

% 番号なしセクション（\section* の代替）
\newcommand{\plainsection}[1]{%
  \vspace{1.5em}%
  \noindent
  \colorbox{sectionblue}{%
    \parbox{\linewidth}{%
      \vspace{2pt}%
      {%
        \Large\bfseries\sffamily\color{black}%
        \JPSans % ← 和文だけ Noto Sans JP に切り替える
        #1%
      }%
      \vspace{2pt}%
    }%
  }%
  \vspace{1.2em}%
}



% --- Subsection：黒ゴシック太字＋下に青ライン ---

\renewcommand{\thesubsection}{\thesection.\arabic{subsection}}

\newcommand{\SubsectionHeading}[1]{%
  \refstepcounter{subsection}%
  %\addcontentsline{toc}{subsection}{\protect\numberline{\thesubsection}#1} % 章番号を目次にも出力する場合
  \addcontentsline{toc}{subsection}{#1} % 章番号を目次にも出力しない場合
  \par\vspace{1.0em}% ← ① 平文との間に確実に改行を入れる
  \noindent % ← ★ これが重要：自動インデントを消す
  {\large\bfseries\sffamily\JPSans\color{black}%
    % \thesubsection\quad #1% ← ★ 章番号を出力する
    #1% ← ★ 章番号を出力しない
  }%
  \par\vspace{-1.2em}% ← ② 見出しとラインの距離を詰める
  \noindent % ← ★ これが重要：自動インデントを消す
  {\color{sectionblue}\rule{\linewidth}{2pt}}%
  %\vspace{0.6em}% ← ラインの下の余白
  \par\vspace{0.4em}% ← ★ 見出しの後に確実に改行＋余白
}

\renewcommand{\subsection}[1]{%
  \SubsectionHeading{#1}%
}

% --- Subsubsection：本文サイズ＋黒ゴシック太字 ---
\newcommand{\SubsubsectionHeading}[1]{%
  \par\vspace{0.8em}% 本文との間隔
  \noindent
  {\normalsize\bfseries\sffamily\JPSans\color{black}%
    #1% ← ★ 章番号を出力しない
  }%
  \par\vspace{0.4em}% ← ★ 見出しの後に確実に改行＋余白
}

\renewcommand{\subsubsection}[1]{%
  \SubsubsectionHeading{#1}%
}

% --- 見出しつきの箇条書きを太文字にする設定 
\renewcommand*\descriptionlabel[1]{%
  \hspace\labelsep
  {\bfseries\sffamily\JPSans\color{black} #1}%
}
% ----------------------------------------
% コラム枠表示（背景色・枠線・フォント調整）
% ----------------------------------------


\newtcolorbox{columnbox}[1][]{
  enhanced,
  breakable,
  colback=white,
  colframe=framegray,
  colbacktitle=framegray,
  coltitle=black,
  fonttitle=\bfseries\sffamily\JPSans,
  title={#1},
  fontupper=\sffamily\JPSans\footnotesize,
  %before upper={\setlength{\parindent}{1em}},
  before upper={
    \addcontentsline{toc}{subsection}{#1}%
    \setlength{\parindent}{1em}%
  },
 before skip=2.0em,
  after skip=2.0em,
  left=8pt,
  right=8pt,
  top=8pt,
  bottom=8pt,
  boxrule=0.4pt
}
% ============================================================
% code-style.tex
% A5 技術書テンプレート用：listings + tcolorbox のコード環境
% ============================================================
% listing の枠線・背景を消す
%	\floatstyle{plain}
%	\restylefloat{listing}

% listing の上下余白を詰める
%	\setlength{\intextsep}{0pt}
%	\setlength{\textfloatsep}{0pt}

% ------------------------------------------------------------
% 必要パッケージ
% ------------------------------------------------------------
% \usepackage{listingsutf8}   % 日本語対応 listings
% \usepackage[most]{tcolorbox}
% \usepackage{xcolor}
\lstset{inputencoding=utf8} 

% ------------------------------------------------------------
% 色定義（コード用）
% ------------------------------------------------------------
\definecolor{codebg}{HTML}{F5F5F5}
\definecolor{keywordblue}{HTML}{005CC5}
\definecolor{stringred}{HTML}{D73A49}
\definecolor{commentgray}{HTML}{777777} % ← 行番号・コメントの色
\definecolor{framegray}{HTML}{D0D7DE}

% ------------------------------------------------------------
% listings 共通スタイル
% ------------------------------------------------------------
\lstdefinestyle{mypython-base}{
  language=Python,
  backgroundcolor=\color{codebg},
  % basicstyle=\codejfont\ttfamily\small,   % ← 和文＋欧文ゴシック体
  basicstyle=\ttfamily\small,
  keywordstyle=\color{keywordblue},
  stringstyle=\color{stringred},
  commentstyle=\codejfont\color{commentgray},
  showstringspaces=false,
  breaklines=true,
  tabsize=4,
  columns=fullflexible,
}

% ------------------------------------------------------------
% 行番号あり
% ------------------------------------------------------------
\lstdefinestyle{mypython}{
  style=mypython-base,
  numbers=left,
  numberstyle=\codejfont\color{commentgray}\small, % ← 太字を外して控えめに
  numbersep=5pt,   % 行番号とコードの距離
}

% ------------------------------------------------------------
% 行番号なし
% ------------------------------------------------------------
\lstdefinestyle{mypython-nonum}{
  style=mypython-base,
  numbers=none,
}

% ------------------------------------------------------------
% tcolorbox によるコードブロック定義
% ------------------------------------------------------------

% 行番号あり：開始行をオプションで指定（デフォルト 1）
\newtcblisting{pycode}[1][1]{%
  listing only,
  listing style=mypython,
  colback=codebg,
  colframe=framegray,
  boxrule=0.4pt,
  arc=3pt,
  outer arc=3pt,
  left=2.0em,   % ← 行番号あり／なしで左端を揃える
  right=4pt,
  top=2pt,
  bottom=2pt,
  enhanced,
  breakable,
  before upper=\lstset{firstnumber=#1},
}

% 行番号なし
\newtcblisting{pycodenonum}{%
  listing only,
  listing style=mypython-nonum,
  colback=codebg,
  colframe=framegray,
  boxrule=0.4pt,
  arc=3pt,
  outer arc=3pt,
  left=2.0em,   % ← 行番号なしでも左端を揃える
  right=4pt,
  top=4pt,
  bottom=4pt,
  enhanced,
  breakable,
}

% ============================================================
% ここまで
% ============================================================


\title{（使わない）}
\author{（使わない）}	
\date{}%これで日付を非表示にできる

% ================================================================

\begin{document}
%\maketitle

%\chapter*{はじめに}
\frontmatter
\tableofcontents
\mainmatter

\newpage
\chapter{pLaTeX から LuaLaTeX へ}
\label{chap:architecture-basics}
TeXは長年pLaTeXが使われていたが、 LuaLaTeX への移行が進んでいる。本章ではその背景とLuaLaTeX の優位性について説明します。


\newpage
\section{なぜ今 LuaLaTeX なのか}

2024年11月、LaTeXチーム（LaTeX Project）は LaTeX News 40 にて、
LuaTeX 以外のエンジンのフルサポートを段階的に終了する方針を公式に発表しました。

日本国内でも、定番の入門書『美文書作成入門』が第9版（2023年12月）から LuaLaTeX を標準エンジンとして採用しています。
また、大学などの教育機関でも pLaTeX から LuaLaTeX への移行が進んでいます。

\section{LuaLaTeX の優位性}
\subsection{フォントの自由度}

LuaLaTeX では \verb|\setmainjfont、\newjfontfamily| などのコマンドにより、OS にインストールされた任意のフォントをウェイト・スタイル単位で細かく制御できます。
このため、市販書籍レベルのTeX文書の生成が容易に可能となります。
一方、pLaTeXではフォント制御が限定的です。

\vspace{5mm}
\subsection{Unicode 完全対応}
pLaTeX は1980年代、Unicode 1.0 リリース以前に設計されたため、限定された文字集合しか扱えません。
LuaLaTeX ではUTF-8をネイティブ入力とし、Unicode全範囲に対応しています。

\vspace{5mm}
\subsection{PDF 直接出力、PDF/A準拠・メタデータ管理}
pLaTeX は DVI ファイルを中間生成し dvipdfmx で PDF に変換する2段階プロセスが必要ですが、
LuaLaTeX は pdfTeX の後継として PDF を直接出力することができます。
これによりビルド工程が簡素化され、PDF のバージョン指定や PDF/A 準拠もエンジンレベルで制御可能となっています。

また、LuaLaTeX では hyperref + hyperxmp（または pdfx）を用いることで、以下が容易に実現できます。
\begin{description}
	\item [PDF/A-3u 準拠] 長期保存・アーカイブ用のISO規格に対応
	\item [XMP メタデータ] タイトル・著者・キーワード等を標準的な形式で埋め込み
	\item [PDF バージョン指定] pdfversion=1.7 等を明示的に設定可能
	\item [日本語しおり（ブックマーク）] pxjahyper 等の追加パッケージなしで Unicode しおりが正常動作
\end{description}

\vspace{15mm}
\begin{columnbox}[{\faLightbulb\hspace{0.5em} 現時点ではpLaTeX の方が有利な点}]
以下の点については、LuaLaTeX よりもpLaTeX 側の方が利点がある。
	\begin{description}
	\item [コンパイル速度] 
	ベンチマーク比較では、pLaTeX が平均 2.98秒に対し LuaLaTeX は平均 7.68秒と、約2.5倍の差 がある。
	大規模文書ではこの差が顕著になる。特にオンライン環境ではタイムアウトが発生するケースも報告されている。

	\item[レガシー資産との互換性] 日本の学会・学術誌が提供するスタイルファイルの多くは pLaTeX（jsarticle / jsbook ベース）向けに作成されている。
	これらをそのまま使用する場合、pLaTeX が必要となるケースがある。

	\item[メモリ使用量]LuaTeX-ja は日本語フォント処理のためメモリ消費が大きく、一部のフォントでは64ビット処理系が必須となる。
	\end{description}
\end{columnbox}

\chapter{テンプレート構成と使い方}
\newpage
\section{preamble フォルダ構成}

本テンプレートでは、LaTeX のプリアンブルを役割ごとに分割し、
保守性と再利用性を高めています。各ファイルの役割は以下の通りです。

\begin{verbatim}
preamble/
 ├ packages.tex
 ├ fonts.tex
 ├ colors.tex
 ├ layout.tex
 ├ header-footer.tex
 ├ chapter-style.tex
 ├ section-style.tex
 ├ column-style.tex
 └ code-style.tex
\end{verbatim}

\begin{description}
  \item[preamble/packages.tex] 使用するパッケージをまとめて読み込みます。
  日本語処理、色、図表、コード、見出し、余白、ヘッダ・フッタなど、
  テンプレート全体の基盤となる設定が含まれます。

  \item[preamble/fonts.tex] 欧文・和文フォントの設定を行います。
  Noto 系フォントを中心に、本文・見出し・コード用フォントを定義します。

  \item[preamble/colors.tex] セクション帯、章扉、コラム、コード枠線など、
  テンプレートで使用する色を定義します。

  \item[preamble/layout.tex] A5 判のページレイアウト、余白、行間などを設定します。
  見開き（twoside）に最適化されています。

  \item[preamble/header-footer.tex] fancyhdr によるヘッダ・フッタ（柱）の設定です。
  左ページに章タイトル、右ページにセクションタイトルを表示します。

  \item[preamble/chapter-style.tex] 章扉のデザインを定義します。
  背景帯、章番号、章タイトルのレイアウトを含みます。

  \item[preamble/section-style.tex] セクション見出しのデザインを定義します。
  青帯＋黒文字＋ゴシック太字のスタイルで統一されています。
  番号なしセクション用の \texttt{\textbackslash plainsection} もここに含まれます。

  \item[preamble/column-style.tex] コラム（囲み記事）のデザインを定義します。
  tcolorbox を用いた枠線・背景色・アイコンなどの設定が含まれます。

  \item[preamble/code-style.tex] listings によるコード表示の設定です。
  背景色、行番号、キーワード色、枠線などを統一しています。
\end{description}

\section{テンプレートの使い方}

\subsection{基本構造}

本テンプレートを使用する際の基本構造は以下の通りです。

\begin{pycode}[1]
\documentclass[10pt,twoside]{book}
% 日本語（LuaLaTeX）
	\usepackage{luatexja}
	\usepackage{luatexja-fontspec}

% 色
	\usepackage{xcolor}

% コード
	\usepackage{listings}
	\usepackage{tcolorbox}
	\tcbuselibrary{skins, breakable, listings}

% 見出し
	\usepackage{titlesec}

% 行間
	\usepackage{setspace}

% 余白
	% \usepackage{geometry}
	% layout.tex で geometry のオプション指定があるので、ここでは読み込まない

% ヘッダ・フッタ
	\usepackage{fancyhdr}

% 図表
	\usepackage{graphicx}
	\usepackage{caption}

% TikZ（章扉）
	\usepackage{tikz}

% 目次
	\usepackage{tocloft}

% コラムアイコン
	\usepackage{fontawesome5}

% listing の位置調整
	\usepackage{placeins}
	\usepackage{float}

% ---------------------------------------------------------
% 日本語フォント設定
% ---------------------------------------------------------
% 日本語フォント（Windows）
\setmainfont{Noto Serif}      % 欧文本文：セリフ
\setsansfont{Noto Sans}       % 欧文ゴシック
\setmainjfont{Noto Serif JP}  % 和文明朝（本文）
\setsansjfont{Noto Sans JP}   % 和文ゴシック（見出し用）

% 日本語の太字を有効化（luatexja）
\ltjsetparameter{jacharrange={-2}} % 日本語の太字を許可
% 日本語ゴシック体（キャプション専用）
\newjfontfamily\captionjfont{Noto Sans JP}[
  UprightFeatures = { Weight = 400 }, % Regular
  BoldFeatures    = { Weight = 700 }  % Bold
]
\newjfontfamily\JPSans{Noto Sans JP}[
  UprightFeatures = { Weight = 400 }, % Regular
  BoldFeatures    = { Weight = 700 }  % Bold
]
% コード用日本語ゴシック体
\newjfontfamily\codejfont{Noto Sans JP}[
  UprightFeatures = { Weight = 400 },
  BoldFeatures    = { Weight = 700 }
]
\definecolor{chapterblue}{HTML}{DCE4F2}

\definecolor{sectionblue}{HTML}{c9d1e8} %よりグレーに近い青に変更（6F8FBC → c9d1e8）
%\definecolor{sectionblue}{HTML}{4A6FA5}

\definecolor{framegray}{HTML}{D0D7DE}

\setstretch{1.4}

\usepackage[
  paperwidth=148mm,
  paperheight=210mm,
  inner=20mm,
  outer=10mm,
  top=18mm,
  bottom=20mm,
  heightrounded,
  twoside
]{geometry}

% ---------------------------------------------------------
% ヘッダ・フッタ（落ち着いた市販書籍風）
% ---------------------------------------------------------
\pagestyle{fancy}

% --- ヘッダのフォント（ゴシック体） ---
\newcommand{\HeaderFont}{\small\sffamily\JPSans}

% --- 章タイトル（chapter）を記録 ---
\renewcommand{\chaptermark}[1]{%
  \markboth{\thechapter\quad #1}{}%
}


% --- 節タイトル（section）を記録 ---
\renewcommand{\sectionmark}[1]{%
  \markright{\thesection\quad #1}%
}


% --- ヘッダ設定 ---
\fancyhf{} % 既存の設定をクリア

% 偶数ページ（左ページ）：章タイトル
\fancyhead[LE]{\HeaderFont \leftmark}

% 奇数ページ（右ページ）：節タイトル
\fancyhead[RO]{\HeaderFont \rightmark}

% ページ番号を外側寄せ
\fancyfoot[LE]{\HeaderFont \thepage} % 偶数ページ（Left Even）
\fancyfoot[RO]{\HeaderFont \thepage} % 奇数ページ（Right Odd）

% ヘッダの下線を消す（技評風）
\renewcommand{\headrulewidth}{0pt}

% --- plain スタイル（目次・章扉などで使われる）を外側寄せに再定義 ---
\fancypagestyle{plain}{
  \fancyhf{}
  % ページ番号を外側寄せ
  \fancyfoot[LE]{\HeaderFont \thepage} % 偶数ページ（左）
  \fancyfoot[RO]{\HeaderFont \thepage} % 奇数ページ（右）
  \renewcommand{\headrulewidth}{0pt}   % ヘッダ線なし（技評風）
}


% 章扉背景（ページ全面を確実に塗りつぶす）
\newcommand{\ChapterBackground}{%
  \begin{tikzpicture}[remember picture,overlay]
    % ページ全体より少し大きめに塗る（塗り残し防止）
    \fill[chapterblue]
      ([xshift=-5mm,yshift=5mm]current page.north west)
        rectangle
      ([xshift=5mm,yshift=-5mm]current page.south east);
  \end{tikzpicture}%
}

% 章デザイン（bbe風：背景色＋装飾）
\titleformat{\chapter}[display]
  {\bfseries\Huge}
  {%
    \ChapterBackground
    \vspace*{3cm}
    \filright\Large\textsf{CHAPTER\ \thechapter}
  }
  {1ex}
  {%
    \titlerule[1.2pt]
    \vspace{2ex}
    \filright\Huge
  }
  [\vspace{2cm}]
% --------------------------------------------------------

% セクション見出しデザイン
% --------------------------------------------------------
% 見出しデザイン（技術書風）


\makeatletter
\def\SectionTitle{}
\def\ttl@save@title#1{\gdef\SectionTitle{#1}}
\makeatother

% --- Section：青背景＋黒文字＋ゴシック太字 ---
\let\oldsection\section

\renewcommand{\section}[1]{%
  \refstepcounter{section}%
  \SectionHeading{\thesection}{#1}%
}

\newcommand{\SectionHeading}[2]{%
  \sectionmark{#2} % ← ★ これが右ページの柱に必要
  \addcontentsline{toc}{section}{\protect\numberline{\thesection}#2}
  \vspace{1.5em}%
  \noindent
  \colorbox{sectionblue}{% sectionblueから変更
    \parbox{\linewidth}{%
      \vspace{2pt}%
      {%
        \Large\bfseries\sffamily\color{black}%
        \JPSans % ← 和文だけ Noto Sans JP に切り替える
        #1\quad #2%
      }%
      \vspace{2pt}%
    }%
  }%
  \vspace{1.2em}%
}

% 番号なしセクション（\section* の代替）
\newcommand{\plainsection}[1]{%
  \vspace{1.5em}%
  \noindent
  \colorbox{sectionblue}{%
    \parbox{\linewidth}{%
      \vspace{2pt}%
      {%
        \Large\bfseries\sffamily\color{black}%
        \JPSans % ← 和文だけ Noto Sans JP に切り替える
        #1%
      }%
      \vspace{2pt}%
    }%
  }%
  \vspace{1.2em}%
}



% --- Subsection：黒ゴシック太字＋下に青ライン ---

\renewcommand{\thesubsection}{\thesection.\arabic{subsection}}

\newcommand{\SubsectionHeading}[1]{%
  \refstepcounter{subsection}%
  %\addcontentsline{toc}{subsection}{\protect\numberline{\thesubsection}#1} % 章番号を目次にも出力する場合
  \addcontentsline{toc}{subsection}{#1} % 章番号を目次にも出力しない場合
  \par\vspace{1.0em}% ← ① 平文との間に確実に改行を入れる
  \noindent % ← ★ これが重要：自動インデントを消す
  {\large\bfseries\sffamily\JPSans\color{black}%
    % \thesubsection\quad #1% ← ★ 章番号を出力する
    #1% ← ★ 章番号を出力しない
  }%
  \par\vspace{-1.2em}% ← ② 見出しとラインの距離を詰める
  \noindent % ← ★ これが重要：自動インデントを消す
  {\color{sectionblue}\rule{\linewidth}{2pt}}%
  %\vspace{0.6em}% ← ラインの下の余白
  \par\vspace{0.4em}% ← ★ 見出しの後に確実に改行＋余白
}

\renewcommand{\subsection}[1]{%
  \SubsectionHeading{#1}%
}

% --- Subsubsection：本文サイズ＋黒ゴシック太字 ---
\newcommand{\SubsubsectionHeading}[1]{%
  \par\vspace{0.8em}% 本文との間隔
  \noindent
  {\normalsize\bfseries\sffamily\JPSans\color{black}%
    #1% ← ★ 章番号を出力しない
  }%
  \par\vspace{0.4em}% ← ★ 見出しの後に確実に改行＋余白
}

\renewcommand{\subsubsection}[1]{%
  \SubsubsectionHeading{#1}%
}

% --- 見出しつきの箇条書きを太文字にする設定 
\renewcommand*\descriptionlabel[1]{%
  \hspace\labelsep
  {\bfseries\sffamily\JPSans\color{black} #1}%
}
% ----------------------------------------
% コラム枠表示（背景色・枠線・フォント調整）
% ----------------------------------------


\newtcolorbox{columnbox}[1][]{
  enhanced,
  breakable,
  colback=white,
  colframe=framegray,
  colbacktitle=framegray,
  coltitle=black,
  fonttitle=\bfseries\sffamily\JPSans,
  title={#1},
  fontupper=\sffamily\JPSans\footnotesize,
  %before upper={\setlength{\parindent}{1em}},
  before upper={
    \addcontentsline{toc}{subsection}{#1}%
    \setlength{\parindent}{1em}%
  },
 before skip=2.0em,
  after skip=2.0em,
  left=8pt,
  right=8pt,
  top=8pt,
  bottom=8pt,
  boxrule=0.4pt
}
% ============================================================
% code-style.tex
% A5 技術書テンプレート用：listings + tcolorbox のコード環境
% ============================================================
% listing の枠線・背景を消す
%	\floatstyle{plain}
%	\restylefloat{listing}

% listing の上下余白を詰める
%	\setlength{\intextsep}{0pt}
%	\setlength{\textfloatsep}{0pt}

% ------------------------------------------------------------
% 必要パッケージ
% ------------------------------------------------------------
% \usepackage{listingsutf8}   % 日本語対応 listings
% \usepackage[most]{tcolorbox}
% \usepackage{xcolor}
\lstset{inputencoding=utf8} 

% ------------------------------------------------------------
% 色定義（コード用）
% ------------------------------------------------------------
\definecolor{codebg}{HTML}{F5F5F5}
\definecolor{keywordblue}{HTML}{005CC5}
\definecolor{stringred}{HTML}{D73A49}
\definecolor{commentgray}{HTML}{777777} % ← 行番号・コメントの色
\definecolor{framegray}{HTML}{D0D7DE}

% ------------------------------------------------------------
% listings 共通スタイル
% ------------------------------------------------------------
\lstdefinestyle{mypython-base}{
  language=Python,
  backgroundcolor=\color{codebg},
  % basicstyle=\codejfont\ttfamily\small,   % ← 和文＋欧文ゴシック体
  basicstyle=\ttfamily\small,
  keywordstyle=\color{keywordblue},
  stringstyle=\color{stringred},
  commentstyle=\codejfont\color{commentgray},
  showstringspaces=false,
  breaklines=true,
  tabsize=4,
  columns=fullflexible,
}

% ------------------------------------------------------------
% 行番号あり
% ------------------------------------------------------------
\lstdefinestyle{mypython}{
  style=mypython-base,
  numbers=left,
  numberstyle=\codejfont\color{commentgray}\small, % ← 太字を外して控えめに
  numbersep=5pt,   % 行番号とコードの距離
}

% ------------------------------------------------------------
% 行番号なし
% ------------------------------------------------------------
\lstdefinestyle{mypython-nonum}{
  style=mypython-base,
  numbers=none,
}

% ------------------------------------------------------------
% tcolorbox によるコードブロック定義
% ------------------------------------------------------------

% 行番号あり：開始行をオプションで指定（デフォルト 1）
\newtcblisting{pycode}[1][1]{%
  listing only,
  listing style=mypython,
  colback=codebg,
  colframe=framegray,
  boxrule=0.4pt,
  arc=3pt,
  outer arc=3pt,
  left=2.0em,   % ← 行番号あり／なしで左端を揃える
  right=4pt,
  top=2pt,
  bottom=2pt,
  enhanced,
  breakable,
  before upper=\lstset{firstnumber=#1},
}

% 行番号なし
\newtcblisting{pycodenonum}{%
  listing only,
  listing style=mypython-nonum,
  colback=codebg,
  colframe=framegray,
  boxrule=0.4pt,
  arc=3pt,
  outer arc=3pt,
  left=2.0em,   % ← 行番号なしでも左端を揃える
  right=4pt,
  top=4pt,
  bottom=4pt,
  enhanced,
  breakable,
}

% ============================================================
% ここまで
% ============================================================


\begin{document}

\tableofcontents

\chapter{はじめに}
本文…

\section{セクションタイトル}
本文…

\plainsection{番号なしセクション}
本文…

\end{document}
\end{pycode}

\subsection{番号なしセクションについて}

標準の \texttt{\textbackslash section*} は本テンプレートのデザインと互換性がないため、
番号なしセクションには \texttt{\textbackslash plainsection} を使用します。


\subsection{PDF ビューワの見開き設定}

Adobe Reader を使用する場合は、以下の設定を推奨します。

\begin{itemize}
  \item 見開きページ表示
  \item 表紙を見開きに表示
\end{itemize}

これにより、左ページ＝偶数、右ページ＝奇数の正しい製本順で表示されます。

\subsection{表紙について}

本テンプレートでは \texttt{\textbackslash maketitle} を使用しません。
表紙は後で別 PDF として作成する方針です。

\section{注意点}

\begin{itemize}
  \item \texttt{\textbackslash section*} は使用しないでください。
  \item コードは listings を使用し、tcolorbox で囲む構造です。
  \item 章扉は自動的にページ番号なしになります。
  \item レイアウトは A5 判に最適化されています。
\end{itemize}

\end{document}
